\section{Scope, limitations, and further research}
\label{sec:questions}

The uniform result is stronger than taking countably many individually
avoided ratios, but it does not remove all restrictions on an infinite
pattern. The next proposition makes one limitation precise.

\begin{proposition}[Ratio bounds alone cannot support uniform avoidance]
\label{prop:annular-obstruction}
Every measurable $E\subset\R$ of positive measure contains a translate
of some positive decreasing sequence $(a_n)$ with
\begin{equation}\label{eq:annular-ratios}
 \frac18<\frac{a_{n+1}}{a_n}<\frac12\qquad(n\ge1).
\end{equation}
Consequently no positive-measure set can simultaneously avoid all
sequences satisfying even these fixed two-sided ratio bounds.
\end{proposition}

\begin{proof}
Choose a Lebesgue density point $x$ of $E$. For all sufficiently large
$j$, the interval
\[
 (x+\tfrac12 4^{-j},\ x+4^{-j})
\]
meets $E$: otherwise its missing length $\tfrac12 4^{-j}$ would
contradict density one at $x$ on intervals of radius $4^{-j}$.
Choose $y_j$ in each such intersection and put $a_j=y_j-x$.
The bounds $\tfrac12 4^{-j}<a_j<4^{-j}$ give
\eqref{eq:annular-ratios}; reindex the sufficiently late tail from one.
\end{proof}

The polynomial dependence in \cref{thm:poly-family} prevents this
unrestricted freedom to choose each annular point independently. The
number of parameters and the growth of their sign complexity are doing
real work. A common envelope is necessary for the present proof but is
not, by itself, enough.

\subsection{What the theorems do not claim}

The general Erd\H{o}s similarity conjecture for an arbitrary infinite
set is not resolved here. In particular, arbitrary rapidly decaying
sequences, such as sequences whose logarithmic scales have unbounded
successive gaps, need not admit the synchronized sampling used above.
The same applies to arbitrary infinite mixtures and to unrestricted
nonlinear images of geometric sequences.

No effective complexity bound of practical size is supplied. The first
branching choice already grows with $\log(1/Q)/p$, while the tree depth
needed to make default avoidance rare can be exponential in the
branching number. The proof explicitly permits this cost. The exact
verification suite tests structural identities at feasible sizes; it
does not instantiate the full small-density construction.

The inherited routing argument is fully credited to \cite{OpenAI84}.
The novel claims relative to that inspected source are the common-ratio
uniformization, the polynomial-family extension and its applications,
the periodic tail assembly, and the effective certificate framework.
A targeted repository and literature search located no earlier statement
with these simultaneous quantifiers. That search is evidence about
provenance, not a substitute for external mathematical review.

\subsection{Questions arising from this work}

The following are proposed research directions. Except for the explicitly
identified classical conjecture, they are questions suggested by this
paper, not assertions that each has an established open-problem status
in the literature.

\begin{question}[Quantitative normalized blockers]
For a ratio interval $[\alpha,\beta]$ and hole budget $\delta$, how small
can the number of rational arcs and the finite witnessing exponent be
in \cref{prop:effective-H}? Can the tree depth be replaced by a
substantially smaller construction?
\end{question}
The present proof controls the parameter entropy locally but pays a
large price for finding a default on every center route with high
probability. A different routing rule could change that price without
changing the synchronization lemma.

\begin{question}[Derandomization with short certificates]
Can one construct useful normalized blockers by deterministic
conditional expectation, while representing the tables implicitly?
Can the final quantified polynomial certificate be made small enough
for complete independent checking in Lean or another proof assistant?
\end{question}
The terminal-address argument gives an exact probability identity.
The obstacle is the number of table entries, not a need to estimate
that identity numerically.

\begin{question}[Infinite mixtures]
Is there one positive-measure compact set avoiding every affine
sequence $a_n=\int_{[\alpha,\beta]}q^n\,d\mu(q)$, where $\mu$ ranges
over all probability measures on a fixed compact interval
$0<\alpha<\beta<1$?
\end{question}
Finite support is covered here with no bound on its size, using a
countable union of finite-dimensional families. Passing to limits of
measures does not preserve a useful uniform finite witness automatically.
The unrestricted-annulus obstruction does not directly settle this
more structured class.

\begin{question}[Bounds near cancellations and root collisions]
For the signed exponential polynomials in \cref{thm:signed-exp-poly},
can one obtain sharp uniform bounds on the first usable tail and the
finite certificate size as the dominant coefficient tends to zero or
as the largest two bases coalesce?
\end{question}
The existence theorem uses a countable exhaustion by coefficient and
root-separation bounds. It gives an effective answer on each compact
box, but does not control the deterioration sharply near its boundary.

\begin{question}[Oscillatory recurrences]
What remains true for real linear recurrence sequences whose
characteristic roots are nonreal and have modulus less than one?
Which spectral separation conditions guarantee a uniformly usable
subsequence and a polynomial parameter count?
\end{question}
Oscillation can create arbitrarily small terms and destroy monotonicity.
Using several coordinates of a recurrence state rather than one scalar
term may provide a suitable replacement for the first-crossing rule.

\begin{question}[A parameter-complexity threshold]
What is the weakest quantitative hypothesis on the number of realized
local key vectors that ensures a common avoiding set? Is there a
natural sharp boundary between the finite-dimensional families handled
here and the class excluded by \cref{prop:annular-obstruction}?
\end{question}
The proof can absorb any fixed exponential rate in window length by
increasing $M$. It cannot absorb arbitrary rates whose coefficient
grows uncontrollably with the window location.

\begin{question}[Sharper effective representations]
Can the compact avoiding set be chosen computable in the Hausdorff
representation, or with an explicitly computable distance function,
while retaining simultaneous avoidance and near-full measure?
\end{question}
The effective measure modulus in \eqref{eq:measure-modulus} does not
by itself answer this. Nor does it automatically make the perfect
support construction effective.

\begin{question}[Geometric regularity of the avoiding set]
Can one prescribe additional regularity of the complementary intervals,
such as quantitative bounds on their distribution, while avoiding all
geometric progressions? Which familiar positive-measure Cantor
constructions necessarily contain some such progression?
\end{question}
The results of \cite{BGKMW} concerning symmetric Cantor sets show
that geometric regularity can force progressions. A successful avoiding
construction must depart from those sufficient conditions.

\begin{question}[Beyond contracting geometric scales]
Can the routing mechanism be combined with a different sampling rule
to handle a fixed supergeometric sequence, or a nontrivial compact
family of them? Can this yield new cases of the full Erd\H{o}s
similarity conjecture?
\end{question}
The obstruction is visible in the proof: a large gap between successive
logarithmic scales leaves too few independent samples at a common
resolution. A new ingredient is required there.

\begin{question}[Uniform nonlinear images]
For which finite-dimensional classes of maps $\phi$ can one avoid all
sets $\{x+s\phi(q^n):n\ge1\}$ simultaneously? Every nonconstant real
polynomial $\phi$ is covered by \cref{thm:signed-exp-poly}, after
absorbing its constant term into the translation;
which additional analytic classes have comparably controlled local
sign complexity?
\end{question}
A Taylor expansion alone does not justify this extension, since a
small error can move a point across the boundary of the avoiding set.
The finite positive-margin certificates provide a possible way to
quantify such errors on compact parameter ranges.
